\section{The algebra and a uniform additivity defect}
\label{sec:algebra}

We now put all the bundle tests into one algebra.  The uniform number of
maps and their uniform norm bounds allow this construction, while each
test obstructs polynomial candidates at its prescribed matrix size and
degree.

\begin{theorem}
\label{thm:construction}
There is a separable unital complex $C^*$-algebra
$B=C^*(1,x_1,\ldots,x_{24})$ such that
\begin{equation}
  \sum_{i=1}^{24}x_i^*x_i=1,
  \qquad
  \sum_{i=1}^{24}x_ix_i^*\le\frac23\,1,
  \label{eq:algebra-relations}
\end{equation}
and $M_m(B)$ is not properly infinite for every integer $m\ge1$.
\end{theorem}

\begin{proof}
For each pair of integers $m,R\ge1$, choose the data supplied by
Proposition~\ref{prop:bundle-test}, and denote its base by $Y_{m,R}$.
At a point $y\in Y_{m,R}$ form the Hilbert space
\[
  \mathcal H_{m,R,y}=\bigoplus_{j=0}^{\infty}E_j(y).
\]
For $1\leq i\leq24$, let $S_i^{m,R}(y)$ be the forward shift on this
space whose block from $E_j(y)$ to $E_{j+1}(y)$ is $X_i^{(j)}(y)$.
The column identity in
Proposition~\ref{prop:bundle-test} gives
\[
  X_i^{(j)*}X_i^{(j)}\le I_{E_j},
\]
so every block, and hence every shift, has norm at most one.  These
bounds hold uniformly over all coordinates.

Let $\mathcal D$ be the bounded $C^*$-product of the algebras
$\mathcal B(\mathcal H_{m,R,y})$ over all these indices: its elements are
the operator families of finite supremum norm.  The uniform bound
therefore defines a single tuple in $\mathcal D$,
\[
  x_i=\bigl(S_i^{m,R}(y)\bigr)_{m,R\geq1,\;y\in Y_{m,R}}
  \qquad(1\leq i\leq24).
\]
The column sum of the shifts is the identity on every level.  Their
row sum is zero on level zero and, on level $j\ge1$, is
\[
  \sum_{i=1}^{24}X_i^{(j-1)}X_i^{(j-1)*}
  \le\frac23 I_{E_j}.
\]
Thus \eqref{eq:algebra-relations} holds in $\mathcal D$.  Take
$B=C^*(1,x_1,\ldots,x_{24})\subseteq\mathcal D$.  This is a nonzero
unital algebra, and it is separable: the $*$-polynomials in these
generators with rational complex coefficients form a countable dense
subset.

Fix $m\ge1$ and suppose, towards a contradiction, that $M_m(B)$
contains isometries $s_1,s_2$ with orthogonal ranges.  Their rectangular
concatenation satisfies
\[
  W=[s_1\ s_2]\in M_{m,2m}(B),
  \qquad W^*W=I_{2m}.
\]
By norm density choose $P\in M_{m,2m}(B)$ with $*$-polynomial entries
such that $\|P-W\|<1/5$.  Since $\|W\|=1$,
\[
  \|P^*P-I_{2m}\|
  \le 2\|P-W\|+\|P-W\|^2<\frac12,
\]
and consequently
\begin{equation}
  P^*P\ge\frac12 I_{2m}.
  \label{eq:polynomial-lower-bound}
\end{equation}
After fixing this approximation, choose $R\ge1$ at least the length of
every word occurring in its entries.  We use the coordinate family
belonging to the already constructed test $(m,R)$.

A word of length at most $R$, applied to level zero, takes values in
levels zero through $R$.  Each generator raises the level by one, each
adjoint lowers it by one, and an adjoint applied at level zero gives
zero.  Constant terms preserve level zero.  Restricting the evaluated
matrix $P$ to the input at level zero therefore gives a fiberwise map
\begin{equation}
  P_0:E_0^{\oplus2m}\longrightarrow
  F_{m,R}:=\left(\bigoplus_{j=0}^{R}E_j\right)^{\oplus m}
  \quad\text{over }Y_{m,R}.
  \label{eq:polynomial-bundle-map}
\end{equation}
This is a continuous bundle map.  Indeed, each block of the
restriction is a finite sum of compositions of the continuous maps
$X_i^{(j)}$ and their adjoints, with $j<R$.  This description is
intrinsic to the bundles; no simultaneous trivialization of the
Hilbert fibers or continuity between different tests is required.
The continuation beyond level $R$ does not enter these compositions. Figure~\ref{fig:finite-test} illustrates
this finite propagation and the resulting bundle map.

\begin{figure}[htbp]
\centering
\begin{tikzpicture}[>=Latex, font=\small,
  level/.style={draw, rounded corners=2pt, minimum width=12mm,
                 minimum height=9mm, fill=blue!4}]
  \node[level] (e0) at (0,0) {$E_0$};
  \node[level] (e1) at (1.8,0) {$E_1$};
  \node (dots) at (3.45,0) {$\cdots$};
  \node[level] (er) at (5.1,0) {$E_R$};
  \node[level,fill=gray!8] (beyond) at (7.2,0) {$E_{R+1}$};
  \draw[->] (e0.north east) to[bend left=15] node[above] {$x_i$} (e1.north west);
  \draw[->] (e1.south west) to[bend left=15] node[below] {$x_i^*$} (e0.south east);
  \draw[<->] (e1) -- (dots);
  \draw[<->] (dots) -- (er);
  \draw[<->,dashed] (er) -- (beyond);
  \begin{scope}[on background layer]
    \node[draw=blue!50!black,dashed,rounded corners=3pt,
          inner xsep=4mm,inner ysep=8mm,fit=(e0)(er)] (window) {};
  \end{scope}
  \node[anchor=south,align=center] at (2.55,1.5)
    {A word of length at most $R$, starting at level zero,\\ stays in these levels};
  \node[anchor=north,align=center] at (2.55,-1.5)
    {$P_0:E_0^{\oplus 2m}\longrightarrow
       \bigl(E_0\oplus\cdots\oplus E_R\bigr)^{\oplus m}$};
\end{tikzpicture}
\caption{The test indexed by $(m,R)$, shown schematically. A polynomial
matrix $P$ whose entries contain only words of length at most $R$ maps
its trivial level-zero input into the displayed target. If
$P^*P\ge\tfrac12 I$, this restriction would embed a trivial bundle of
rank $2md_0$ into that target,
contradicting its retained Chern class. The arrows show level changes;
each forward block is one of the maps $X_i^{(j)}$. The continuation after
level $R$ does not affect this test. When $R=1$, the labels $E_1$ and
$E_R$ denote the same level.}
\label{fig:finite-test}
\end{figure}


At each $y$, let $\iota_0$ include the level-zero input in
$\mathcal H_{m,R,y}^{\oplus2m}$, and let $\iota_F$ include
$F_{m,R}(y)$ in $\mathcal H_{m,R,y}^{\oplus m}$.  The range statement
above says $P(y)\iota_0=\iota_F P_0(y)$.  Evaluation preserves
\eqref{eq:polynomial-lower-bound}, so compression of this quadratic-form
inequality gives
\[
  P_0(y)^*P_0(y)
  =\iota_0^*P(y)^*P(y)\iota_0
  \geq\frac12 I
  \qquad(y\in Y_{m,R}).
\]
Hence
\[
  Q=P_0(P_0^*P_0)^{-1/2}
\]
is a continuous fiberwise isometry.  Its image is a subbundle of
$F_{m,R}$ isomorphic to the trivial bundle $E_0^{\oplus2m}$, of rank
$2m\operatorname{rank}(E_0)$.  This contradicts
Proposition~\ref{prop:bundle-test}.  Since $m$ was arbitrary, the
conclusion holds for every matrix algebra over $B$.
\end{proof}

The passage from the two relations to a specified positive pair is
useful independently of the construction.  It requires only the
$1$-quasitrace axioms once a quasitrace on the matrix algebra is given.

\Needspace{12\baselineskip}
\begin{lemma}
\label{lem:defect}
Let $B$ be a unital complex $C^*$-algebra, let $N\ge1$, and suppose
$x_1,\ldots,x_N\in B$ satisfy
\[
  \sum_{i=1}^N x_i^*x_i=1,
  \qquad
  \sum_{i=1}^N x_ix_i^*\le c1
  \quad\text{for some }0\le c<1.
\]
There are positive contractions $a,b\in M_{2N}(B)$, determined by this
tuple, such that every normalized $1$-quasitrace $\tau$ on $M_{2N}(B)$
satisfies
\begin{equation}
  \tau(a+b)-\tau(a)-\tau(b)\ge\frac{1-c}{2N}.
  \label{eq:uniform-defect}
\end{equation}
\end{lemma}

\begin{proof}
Put
\[
  u_i=x_i^*x_i,\qquad v_i=x_ix_i^*,\qquad
  H_i=\sum_{j=1}^i u_j,\qquad J_i=\sum_{j=1}^i v_j,
\]
with $H_0=J_0=0$.  We use these partial sums to define two diagonal
matrices.  In their additivity defect, the first $N$ blocks will telescope
the column partial sums, while the last $N$ use complements to telescope
the row partial sums with the opposite sign and keep the entries positive.
\begin{equation}
\begin{aligned}
  a&=\operatorname{diag}
      (H_0,\ldots,H_{N-1},\,1-J_0,\ldots,1-J_{N-1}),\\
  b&=\operatorname{diag}
      (u_1,\ldots,u_N,\,1-v_1,\ldots,1-v_N).
\end{aligned}
\label{eq:defect-pair}
\end{equation}
All entries are positive contractions: $0\le H_i\le H_N=1$,
$0\le u_i\le1$, and $0\le J_i,v_i\le J_N\le c1$.

Fix a normalized $1$-quasitrace $\tau$ on $M_{2N}(B)$.  Write
$\iota_k(z)$ for $z\in B$ placed in the $k$th diagonal corner and
define $\varphi(z)=\tau(\iota_1(z))$.  If $z\ge0$, let $t$ be the
matrix with $z^{1/2}$ in its $(k,\ell)$ entry and zeros elsewhere.
Then $t^*t=\iota_\ell(z)$ and $tt^*=\iota_k(z)$, so the identity
$\tau(t^*t)=\tau(tt^*)$ gives
\[
  \tau(\iota_k(z))=\varphi(z)
  \qquad (z\ge0, 1\le k\le2N).
\]
Since self-adjoint elements supported in distinct diagonal corners
commute, abelian linearity yields
\begin{equation}
  \tau\bigl(\operatorname{diag}(z_1,\ldots,z_{2N})\bigr)
    =\sum_{k=1}^{2N}\varphi(z_k)
  \quad (z_k\ge0),
  \qquad \varphi(1)=\frac1{2N}.
  \label{eq:corner-normalization}
\end{equation}
The last equality follows by taking every $z_k=1$.  The square axiom
applied within the first corner also gives
\begin{equation}
  \varphi(u_i)=\varphi(v_i)\qquad (1\le i\le N).
  \label{eq:corner-square-equality}
\end{equation}

The first $N$ diagonal blocks contribute to the additivity defect the
terms
\[
  \varphi(H_i)-\varphi(H_{i-1})-\varphi(u_i).
\]
The last $N$ contribute
\begin{align*}
  &\varphi(2-J_i)-\varphi(1-J_{i-1})-\varphi(1-v_i)\\
  &\hspace{2em}=-\varphi(J_i)+\varphi(J_{i-1})+\varphi(v_i).
\end{align*}
Here each scalar-translation identity is taken within the abelian
algebra generated by the corner identity and the indicated
self-adjoint element.  Thus these equalities use no additivity on
noncommuting elements of $B$.  Summing the scalar defects,
telescoping, and using \eqref{eq:corner-square-equality}, we obtain
\begin{align*}
  \tau(a+b)-\tau(a)-\tau(b)
    &=\varphi(H_N)-\sum_{i=1}^N\varphi(u_i)
      -\varphi(J_N)+\sum_{i=1}^N\varphi(v_i)\\
    &=\varphi(1)-\varphi(J_N).
\end{align*}
The restriction of $\varphi$ to $C^*(1,J_N)\subseteq B$ is positive
and linear.  Since $0\le J_N\le c1$, it follows that
$\varphi(J_N)\le c\varphi(1)$.  Equation~\eqref{eq:corner-normalization}
now gives \eqref{eq:uniform-defect}.
\end{proof}
